\begin{center}
    {\Large\itshape Alban au XV\textsuperscript{ème} siècle (suite)}
\end{center}

\bigskip

Ainsi, Saint Martin de Négremont finira par céder le pas à Notre Dame d'Alban qui devient successivement « chef lieu de Claverie » (une ferme porte encore sur la paroisse de Négremont ce nom de « La Claverie »), puis chef lieu de district et enfin de Doyenné.

On s'aperçoit que quelques siècles plus tard, les revenus de ce groupe de paroisses ne sont pas négligeables. En effet, le 13 juin 1721, Monseigneur de la Croix de Castries, évêque d'Albi, donne la jouissance des fruits décimaux de la claverie d'Alban -- et cela pour six ans -- à Maître Jean Dizarn, sieur de la Balme (habitant Coupiac) pour une redevance de 6000 livres : 920 en sont perçues à Saint Pierre d'Illergue, 1721 et 5 sols à Massals et Paulin, 854 à Massuèyes, 660 à Bonneval, 1450 à Teillet, 700 à Saint Jean du Salès. Quant aux « fruits décimaux » de Saint Jean de Janes, c'est le Grand Séminaire d'Albi qui en bénéficie.

À l'heure actuelle, le Doyenné Alban-Valence-Massals-Trébas-Teillet est beaucoup plus étendu géographiquement, mais ne comporte aucun revenu pour le titulaire (dommage !!!). Monsieur le Doyen de N.D. d'Alban.
